\documentclass[11pt,a4paper]{ctexart}

\usepackage[margin=1in]{geometry}
\usepackage{authblk}
\usepackage{booktabs}
\usepackage{array}
\usepackage{hyperref}
\usepackage{graphicx}
\usepackage{amsmath}
\usepackage{xcolor}

\title{Conventional-SC-Dataset：面向超导体系文献与物性参数的可审核数据平台}

\author[1]{[Yuan Ma, Qiang Guo, Xin Zhong, Hanyu Liu, Yanming Ma]}
\affil[1]{[Your Affiliation]}
\date{}

\begin{document}

\begin{document}
\maketitle

\begin{abstract}
超导体研究越来越依赖对分散信息的结构化访问：DOI 连接的来源论文
\textcolor{red}{（DOI连接的来源论文给人感觉不专业，一个专业的数据库不应该专注于DOI来源论文，虽然你的数据源是DOI来源论文，但你应该强调你是一个结构化数据库，提供了丰富的超导体相关信息，而不仅仅是DOI来源论文的连接。建议改为：结构化访问超导体相关信息，包括体系名称、晶体结构、对称性、压力条件、电子结构、电子–声子相关参数以及超导临界转变温度。）}
、体系名称、对称性与结构描述、临界转变温度、压力条件以及电子--声子相关参数。公共材料数据库和近年的超导数据集已经推动了结构--性质分析、基准构建和高通量筛选，但面向具体超导体系的实用记录仍常常保存在静态表格、本地笔记或项目脚本中。本文介绍 Conventional-SC-Dataset，一个可审核的网页平台，用于围绕元素组合组织超导体和超导体系，并将 DOI 连接的论文元数据、每篇论文的多组物理参数、结构信息、贡献者身份、受控可视化和可选择数据源连接起来。平台提供周期表入口、三种标准化的超导体检索模式、基于 DOI 连接的指定体系浏览、DOI 驱动的论文提交、超导记录的表格批量导入、JSON 格式超导数据的导入导出、AI 筛选文献数据库模式、HTSC-2025 基准浏览、Alexandria 电声耦合接口，以及连接结构文件和 PDOS 信息的 Tc 预测页面。在当前仓库快照中，代码库包含 15,572 行后端、前端和维护脚本代码，68 个后端路由和 9 个主要页面模板。主元数据库包含 118 个元素且当前未承载论文记录；另一个 AI 筛选元数据库包含 688 篇论文、1,010 个化合物和 2,624 条物理参数记录，其中 1,921 条包含 Tc 字段，2,068 条包含压力字段；随仓库提供的 HTSC-2025 JSON fixture 包含 140 条基准风格记录。本文不声称提出新的超导机制或经过验证的预测模型，而是建立了一个可复现的数据管理层，用于把分散的超导体记录和相邻计算数据集转化为经过审核、以元素为中心、可复用的数据。
\end{abstract}

\noindent\textbf{关键词：} 超导数据库；常规超导体；氢化物；文献整理；物理参数；数据平台；科研软件

\section{引言}

超导体的寻找已经离不开数据基础设施。常规声子介导超导体可以借助较成熟的第一性原理工作流进行分析，而高压氢化物、常压候选体系和基准数据集不断扩展研究者需要比较的化学与结构空间。近期工作体现了这一趋势：高通量计算已经绘制了实验已知化合物中的 Bardeen--Cooper--Schrieffer 超导体图景 \cite{bercx2025prxenergy}；氢化物机器学习研究筛选了大量候选结构或提出了中等压力超导体的可解释描述符 \cite{pires2026ambient,chen2026descriptors}；HTSC-2025 为常压高温超导体的 AI 驱动临界温度预测提供了基准数据集 \cite{han2025htsc}；HPCSD 高压晶体结构数据库进一步说明了极端条件材料研究中可追溯、压力分辨数据基础设施的重要性 \cite{wang2026hpcsd}。

这些资源解决的是超导数据问题中的不同部分。结构数据库和高通量档案重视晶体结构、能量、相稳定性或计算得到的电子--声子信息；基准数据集重视可比较的模型评估；文献抽取工作重视从文本中转换材料性质记录。然而，在模型训练或机理解释之前，研究者常常面对更操作性的问题：对于一个选定的元素体系，已经报道了哪些论文？每篇论文中录入了哪些物理参数？哪些记录经过了人工审核？哪些记录适合进入公开图表？这不仅是数据科学问题，也是协作整理和质量控制问题。

这个瓶颈在超导体系中尤其明显，因为一篇论文可能同时报告多个化合物、压力区间、结构、样品或计算量。把一篇论文压缩成一个 formula--Tc 对会丢失重要上下文；而把所有证据都写在自由文本笔记中，又会增加后续筛选和比较的难度。因此，一个实用平台需要把论文级元数据和数据点级物理参数分开，保留结构与来源可追踪信息，并把质量控制状态放入工作流，而不是把它包装成独立科学创新点。

Conventional-SC-Dataset 正是为这种需求而开发。平台的基本路径是：选择元素，选择数据源，找到对应超导体系，浏览 DOI 连接的论文或计算候选，提交带有物理参数的记录，在审核界面中整理记录，并把筛选后的数据展示为图表和统计。这个设计借鉴了 HPCSD 等数据库论文的基础设施逻辑：核心贡献不是单个预测，而是一个标准化、可追踪的资料库，使分散研究证据更容易比较和复用。二者的科学对象不同：HPCSD 解析压力相关晶体结构和 DFT 重新计算的能量，而 Conventional-SC-Dataset 解析 DOI 连接的超导体记录、元素组合、结构描述、物理参数条目，以及通往相邻超导数据集的适配器。

本文报告 Conventional-SC-Dataset 的当前实现和边界。我们描述其数据模型、用户工作流、系统实现、本地数据状态，以及与现有超导和材料数据库的关系。同时，我们明确说明当前平台尚不支持的内容：它不是全文文献挖掘引擎，不是经过验证的 Tc 预测模型，不是通用社交社区，也还不是具有冻结版本发布的公共基准。AI 筛选数据库用于快照统计时，只作为模式容量和工作流使用的本地证据；除非增加发布策略、同步流程和持久归档，否则不应被写成已经正式发布的数据集。

\section{相关数据资源}

材料信息学已经形成了几类成熟数据库范式。Materials Project \cite{jain2013materialsproject}、OQMD \cite{kirklin2015oqmd}、AFLOW \cite{curtarolo2015aflow} 和 NOMAD \cite{draxl2018nomad} 展示了标准化计算数据如何支持大规模材料发现；ICSD \cite{belsky2002icsd} 和 COD \cite{grazulis2012cod} 提供了重要的晶体学基础。这些资源建立了材料数据应可检索、可复用并连接到明确元数据的基本预期。

超导专门资源沿着互补方向发展。SuperCon 及其机器可读或本体化衍生工作是历史超导性质数据的重要来源 \cite{nimsSupercon,foppiano2023supercon2,ishii2023ontology}。3DSC 数据集把超导转变温度和晶体结构连接起来，弥补了仅有化学式数据集的一个关键不足 \cite{sommer2023threedsc}。HTSC-2025 聚焦近年的常压高温超导体，并明确定位为 AI 驱动 Tc 预测的基准 \cite{han2025htsc}。高通量计算研究则提供了候选列表和电子--声子性质 \cite{bercx2025prxenergy,pires2026ambient}。

HPCSD 与本文特别相关，因为它把数据库明确写成碎片化领域的基础设施 \cite{wang2026hpcsd}。HPCSD 将实验和理论高压结构整合为压力分辨资料库，初始版本包含跨 89 个元素的 77,346 条结构记录。它的论点是：与常压材料数据库相比，高压结构数据长期稀疏且分散。Conventional-SC-Dataset 采用类似的基础设施立场，但面向另一个缺口：超导体系和物理参数记录在进入下游统计或模型构建之前，需要元素中心浏览、DOI 来源连接、结构描述和人工质量控制。

\section{设计原则}

Conventional-SC-Dataset 通过五个设计原则把整理需求转化为系统约束。首先，平台从元素组合开始，因为这是超导研究者自然的查询单位；其次，平台分离论文元数据和物理参数条目，因为一篇论文可报告多个体系或条件；第三，平台记录质量控制状态，因为提交记录并不自动等于可发布数据；第四，平台把预测结果和已审核文献记录区分开，因为当前证据支持的是基础设施贡献，而不是模型验证贡献；第五，平台把外部或计算数据集作为显式数据源模式处理，而不是静默混入已审核本地语料。

\subsection{元素组合作为主入口}

超导研究者经常从化学组成出发思考问题。因此，平台以周期表作为入口，而不是单纯提供自由文本搜索框。被选择的元素会被标准化、排序，并映射到超导体系。系统支持三种检索模式：\texttt{only}、\texttt{combination} 和 \texttt{contains}。\texttt{only} 检索元素集合完全一致的记录；\texttt{combination} 检索所选元素集合的已有子体系；\texttt{contains} 检索包含所选元素的更大体系。这一区分使同一界面既能支持精确二元或三元检索，也能支持更宽泛的元素族探索。

\subsection{显式数据源模式}

化合物页面现在区分四类浏览模式：已审核本地文献数据库、AI 筛选文献数据库、HTSC-2025 和 Alexandria 电声耦合数据。本地和 AI 筛选模式返回 DOI 连接的论文记录，并支持分页和审核相关筛选。HTSC-2025 模式从 JSON fixture 返回基准风格的常压超导候选。Alexandria 模式实现为本地导入电声耦合数据库之上的适配器，提供元素检索、最低 Tc 筛选、稳定性筛选、材料详情、下载和 CIF 风格结构导出接口。这种数据源分离使来源保持可见：文献记录、基准候选和计算 EPC 条目可以从同一元素工作流检索，但不会被当作同一种证据。

\subsection{论文级元数据与物理参数记录分离}

平台把文献记录和物理数据点分开。\texttt{papers} 表保存 DOI、标题、作者、期刊、年份、摘要、贡献者字段、质量控制状态、审核意见和可视化资格。\texttt{paper\_data} 表保存一个或多个与论文和化合物关联的条目，包括文章类型、超导体类型、化学式、晶体结构、Tc、压力、电子--声子耦合参数、对数平均声子频率、费米能级态密度、样品名、数据来源说明和论文内顺序。这种分离保留了一篇论文可能包含多个超导体系或多个压力--Tc 记录的事实。

\subsection{整理作为可审核工作流}

数据录入被视为工作流，而不是一次性上传。普通用户可以注册、验证邮箱、登录并从化合物页面提交论文；整理者可以审核和编辑提交；更高权限的维护者可以管理用户权限。审核状态、审核人、审核时间、审核意见和可视化资格随论文记录一起保存。这些字段是质量控制基础设施，不是本文要突出的独立科学创新点。

\subsection{预测作为结构输入工具}

仓库包含一个面向结构和 PDOS 输入的 Tc 预测页面。该模块读取上传的 CONTCAR 和 PDOS 文件，包括氢投影 PDOS 信息，并返回用于探索的 Tc 相关估计结果。它不替代文献整理，也不在本文中作为经过验证的预测贡献。这个边界很重要，因为当前证据支持平台作为数据基础设施，而不是支持一个经过基准验证的机器学习模型。

\section{系统实现}

系统由前端页面、后端 API、关系型数据对象、数据源适配器和维护脚本组成。因此，本节从可复现整理工作流的角度描述实现，而不是逐个说明界面控件。

\subsection{数据模型}

当前数据模型包含六类核心对象：元素、化合物、用户、论文、物理参数记录和可视化附件。元素保存 118 个化学元素并支撑周期表界面；化合物保存化学式和标准化元素列表；用户保存普通用户、整理者和维护者状态；论文保存文献元数据和质量控制字段；物理参数记录保存每篇论文下的化合物与超导数据；可视化附件记录当前支持来源图片，并可进一步演进为生成式晶体结构示意图。

实现中还包含非整理型数据适配器。AI 筛选论文 API 复用论文和物理参数响应模型，但面向独立元数据库。HTSC-2025 从仓库 JSON 文件加载，并通过检索、统计和详情接口暴露。Alexandria 支持要求在 Alexandria 数据目录下构建本地 SQLite 数据库，并把重型原始字段放在普通文献表之外；API 响应有意返回轻量材料摘要、结构元数据和生成的 CIF 文本，而不是大型力常数或动力学矩阵负载。

\begin{table}[t]
\centering
\caption{核心数据库对象及其在论文论证中的作用。}
\label{tab:model-zh}
\begin{tabular}{>{\raggedright\arraybackslash}p{0.18\linewidth} >{\raggedright\arraybackslash}p{0.34\linewidth} >{\raggedright\arraybackslash}p{0.36\linewidth}}
\toprule
对象 & 主要字段 & 平台作用 \\
\midrule
\texttt{elements} & 元素符号、英文名、中文名、原子序数 & 周期表渲染与元素合法性校验 \\
\texttt{compounds} & 化学式、元素列表、组成、元素 ID、元素比例 & 元素组合检索和体系分组 \\
\texttt{papers} & DOI、标题、作者、期刊、年份、摘要、贡献者字段、审核字段、图表标记 & 文献级元数据与质量控制状态 \\
\texttt{paper\_data} & 文章类型、超导体类型、化学式、结构、Tc、压力、$\lambda$、$\omega_{\log}$、$N(E_F)$、样品说明 & 每篇论文下的多组物理参数记录 \\
\texttt{paper\_images} & 论文 ID、图标签、原图、缩略图、文件大小 & 可视化附件与来源/结构上下文 \\
\texttt{users} & 邮箱、真实姓名、密码哈希、邮箱验证、管理员标记、审批状态 & 认证、提交、整理和权限控制 \\
\bottomrule
\end{tabular}
\end{table}

\subsection{用户工作流}

主要用户工作流从首页开始。用户在周期表中选择元素，选择检索模式，并进入化合物页面。化合物页面根据所选数据源通过后端 API 检索记录。用户可以按质量控制状态、关键词、年份范围和分页浏览本地及 AI 筛选文献记录；Alexandria 增加最低 Tc 和稳定性筛选；HTSC-2025 提供基准候选分页和详情视图。每张本地或 AI 筛选论文卡片可以展示元数据、DOI、作者、文章类型、超导体类型、化学式、晶体结构、摘要和物理参数条目。

认证用户可以从化合物页面提交论文。提交需要 DOI 信息、元素上下文、文章类型、超导体类型和至少一条物理参数记录。当前实现可以保存可视化附件，但从论文贡献表述上更应强调晶体结构可视化方向，而不是把截图上传当作创新点。后端会验证登录状态，解析 JSON 格式的元素和物理参数负载，检查 DOI 元数据，创建或复用化合物，防止同一化合物体系中的 DOI 重复，创建论文，并写入物理参数记录。

\subsection{整理者工作流}

整理者界面支持平台的质量控制侧。整理者可以检查未审核论文、更新审核状态、编辑记录，并决定某篇论文是否进入首页图表。维护者可以批准整理者申请并更新用户权限。这种角色层次是可信数据管理的支撑基础设施，而不是本文摘要中需要强调的创新点。

\subsection{导入、导出与维护}

平台支持面向批量数据的数据操作。这里的批量导入是指从约定格式的表格或文本文件一次性导入多条超导论文和物理参数记录，而不是手动逐条录入 DOI 记录。维护脚本支持数据库初始化、JSON 格式超导数据导出与导入、ID 迁移、超级管理员创建和备份操作。这些脚本构成平台可复现性的一部分：数据库不仅是网页状态，也是可以初始化、导出、导入和维护的数据对象。

\section{当前仓库与数据状态}

我们检查了 2026 年 6 月 8 日将 \texttt{mayuan} 分支同步到 \texttt{origin/master} 之后的仓库快照。按后端、前端和维护脚本统计，代码库包含 15,572 行 Python、JavaScript、HTML、CSS 和 shell 脚本代码；包含 68 个后端路由和 9 个主要页面模板。当前 \texttt{backend/api} 目录包含 9 个非 \texttt{\_\_init\_\_} API 模块。这些数字是实现描述，而不是科学性能指标；但它们说明平台已经超出静态表格或单个原型脚本。

用于投稿时，这些快照描述应在归档版本中重新生成，并与确切统计协议一起报告。可复现包应包括仓库提交或发布标签、代码行统计包含的文件扩展名、FastAPI 装饰器路由统计命令、表~\ref{tab:status-zh} 使用的数据库文件，以及一个基于 fixture 的工作流检查：初始化数据库、导入小规模论文集、导出、重新导入，并比较记录数和关键字段。相同包还应记录必填提交字段、DOI 元数据失败行为、重复 DOI 规则、审核状态转换、可视化资格和计划中的结构可视化路径。

\begin{table}[t]
\centering
\caption{当前仓库快照中的实现与本地数据状态。}
\label{tab:status-zh}
\begin{tabular}{>{\raggedright\arraybackslash}p{0.42\linewidth} r >{\raggedright\arraybackslash}p{0.34\linewidth}}
\toprule
指标 & 数值 & 范围 \\
\midrule
总代码行数 & 15,572 & 后端、前端和维护脚本 \\
后端路由 & 68 & 后端 FastAPI 路由装饰器 \\
主要页面模板 & 9 & 前端模板目录下的 HTML 页面 \\
\texttt{hydride\_literature.db} 中元素数 & 118 & 主元数据库 \\
\texttt{hydride\_literature.db} 中论文数 & 0 & 当前主元数据库快照 \\
\texttt{hydride\_literature\_ai.db} 中化合物数 & 1,010 & AI 筛选元数据库 \\
\texttt{hydride\_literature\_ai.db} 中论文数 & 688 & AI 筛选元数据库 \\
\texttt{hydride\_literature\_ai.db} 中不同 DOI 数 & 429 & 非空 DOI 去重 \\
\texttt{hydride\_literature\_ai.db} 中物理参数记录数 & 2,624 & \texttt{paper\_data} 表 \\
包含 Tc 字段的记录 & 1,921 & 非空 \texttt{paper\_data.tc} \\
包含压力字段的记录 & 2,068 & 非空 \texttt{paper\_data.tc\_press} \\
AI 筛选论文年份范围 & 1909--2025 & 最小和最大论文年份 \\
AI 筛选论文期刊数 & 105 & 非空期刊名去重 \\
HTSC-2025 fixture 记录数 & 140 & 随仓库提供的 JSON 基准风格候选记录 \\
\bottomrule
\end{tabular}
\end{table}

这些数据状态需要谨慎解释。主元数据库 \texttt{hydride\_literature.db} 在当前快照中包含元素字典和一个化合物记录，但没有论文或物理参数记录。独立的 AI 筛选元数据库 \texttt{hydride\_literature\_ai.db} 包含更多文献与物理参数数据。HTSC-2025 fixture 是作为 JSON 数据源提供的，而不是作为已审核本地文献记录。Alexandria 接口已经实现，但需要本地导入 Alexandria SQLite 数据库后才能报告有意义的行数。因此，平台的数据模型和工作流能够支持非平凡文献语料和外部计算数据浏览，但在形成正式公开版本或同步策略之前，论文应区分默认主数据库、AI 筛选数据文件、随仓库提供的基准 fixture 和可选导入的 EPC 数据库。

\section{与近期数据集的比较}

最接近的概念比较不是某一个数据集，而是一组基础设施工作。HPCSD 标准化高压晶体结构和压力分辨能量，包含来自已报道元素相和 CALYPSO 晶体结构预测任务的两条数据流 \cite{wang2026hpcsd}。HTSC-2025 整理常压高温超导体基准，用于 AI 驱动 Tc 预测 \cite{han2025htsc}。PRX Energy 的 BCS 超导体研究使用高通量第一性原理计算，从实验已知化合物中识别潜在声子介导超导体 \cite{bercx2025prxenergy}。氢化物机器学习研究则通过大规模筛选或可解释描述符扩展候选发现 \cite{pires2026ambient,chen2026descriptors}。

\begin{table}[t]
\centering
\caption{Conventional-SC-Dataset 相对于邻近资源的功能定位。}
\label{tab:comparison-zh}
\begin{tabular}{>{\raggedright\arraybackslash}p{0.18\linewidth} >{\raggedright\arraybackslash}p{0.22\linewidth} >{\raggedright\arraybackslash}p{0.22\linewidth} >{\raggedright\arraybackslash}p{0.25\linewidth}}
\toprule
资源类型 & 主要对象 & 主要优势 & 本平台补充的缺口 \\
\midrule
材料数据库 & 计算材料条目、结构、能量、元数据 & 大规模可复用材料数据与互操作性 & 面向超导体的 DOI 连接、体系记录、物理参数行和可视化资格 \\
结构数据库 & 晶体结构 & 经整理或开放的晶体结构参考数据 & 文献级超导参数和结构感知浏览工作流 \\
SuperCon 类资源 & 超导化学式和性质记录 & 历史超导性质覆盖 & 超导体检索模式、每篇论文多条记录、DOI 连接和网页整理流程 \\
3DSC 与 HTSC-2025 & 结构增强或基准超导数据集 & 结构连接和面向 AI 的基准范围 & 在冻结基准构建之前的动态整理工作流 \\
HPCSD & 压力分辨高压结构 & 极端条件材料的可追踪结构资料库 & 文献和物理参数整理，而非 DFT 重新优化的结构基础设施 \\
\bottomrule
\end{tabular}
\end{table}

Conventional-SC-Dataset 有意定位在这些任务之前的数据管道层级。它不像 HPCSD 那样用统一 DFT 协议重新计算所有结构；不像 HTSC-2025 那样定义冻结基准划分，尽管它能够浏览一个 HTSC-2025 风格 fixture；也不声称完成穷尽式高通量电子--声子计算，尽管它现在提供了面向导入 EPC 记录的 Alexandria 适配器。相反，它支持许多下游任务需要的整理层：把超导体系映射到 DOI 连接的论文，为每篇论文记录多组物理参数，保留结构描述，并控制哪些记录适合可视化和复用。

\section{局限与下一步}

当前平台仍有若干局限。数据模型方面，提交论文和注册用户之间目前主要通过贡献者姓名和单位字段表示，而不是稳定的 contributor-user 外键。存储和维护方面，当前使用 BLOB 保存可视化附件，后续需要向生成式晶体结构可视化过渡，同时也缺少 Alembic 这类正式 schema migration 体系。工作流和检索方面，压力和 Tc 字段已经存在于物理参数记录中，但尚未在每一条本地检索路径中成为成熟的一等筛选维度，尽管 Alexandria 和 HTSC 相关模式已经提供了数据源特定筛选。发布和数据源方面，当前还缺少冻结公共数据集版本、持久归档标识、主数据库与 AI 筛选数据库之间的同步策略，以及带有可复现导入元数据的 Alexandria 数据库。社区方面，当前还缺少点赞、评论、弹幕、热度排序和数据库快照管理。

这些局限定义了实际开发路径。下一版本应优先建立稳定的 \texttt{contributor\_user\_id} 关系、压力和 Tc 区间筛选、更清晰的索引策略、后端管理的引用导出，以及正式 schema migration。只有这些基础稳定之后，公共发布版本、数据集卡、基准划分或社区互动功能才更适合进入数据模型。

\section{结论}

Conventional-SC-Dataset 提供了一个可审核、以元素为中心的超导体系和物理参数整理平台。通过结合周期表入口、标准化超导体检索模式、基于 DOI 的论文连接、每篇论文多组物理参数、结构描述、质量控制状态、可视化资格、JSON 格式超导数据导入导出工具、AI 筛选浏览、HTSC-2025 fixture 浏览和 Alexandria EPC 适配器，它把分散记录和相邻计算数据集转化为可维护的数据工作流。其当前贡献是基础设施性的：它支持超导研究中的已审核数据组织，并为更可靠的统计、模型训练语料和基准构建奠定基础。未来工作应聚焦用户--论文关系建模、物理参数筛选、公共数据版本化、迁移和同步策略，然后再提出更强的预测或社区规模部署主张。

\section*{数据与代码可用性}

源代码和本地数据库文件位于项目仓库 \texttt{[repository URL or local project path]}。在投稿前，应补充正式公共数据发布、版本标识和持久归档 DOI。本文统计来自 2026 年 6 月 8 日将 \texttt{mayuan} 分支同步到 \texttt{origin/master} 之后的仓库快照。

\section*{作者贡献}

\textbf{[Your Name]} 设计并实现平台，整理论文范围，并准备文章草稿。其他贡献者角色应在确认协作者、代码贡献者和数据整理者后补充。

\section*{利益冲突}

作者声明不存在利益冲突。\textbf{[投稿前如有需要请修改。]}

\section*{致谢}

作者感谢本文引用的公共超导、材料和高压数据库研究的开发者与作者。\textbf{[投稿前补充基金号、合作者和单位支持。]}

\begin{thebibliography}{99}

\bibitem{jain2013materialsproject}
A. Jain, S. P. Ong, G. Hautier, W. Chen, W. D. Richards, S. Dacek, S. Cholia, D. Gunter, D. Skinner, G. Ceder, and K. A. Persson.
The Materials Project: A materials genome approach to accelerating materials innovation.
\textit{APL Materials} \textbf{1}, 011002 (2013).
doi:10.1063/1.4812323.

\bibitem{kirklin2015oqmd}
S. Kirklin, J. E. Saal, B. Meredig, A. Thompson, J. W. Doak, M. Aykol, S. Ruehl, and C. Wolverton.
The Open Quantum Materials Database (OQMD): assessing the accuracy of DFT formation energies.
\textit{npj Computational Materials} \textbf{1}, 15010 (2015).
doi:10.1038/npjcompumats.2015.10.

\bibitem{curtarolo2015aflow}
S. Curtarolo, W. Setyawan, G. L. W. Hart, M. Jahnatek, R. V. Chepulskii, R. H. Taylor, S. Wang, J. Xue, K. Yang, O. Levy, M. J. Mehl, H. T. Stokes, D. O. Demchenko, and D. Morgan.
AFLOW: The automatic flow framework for materials discovery.
\textit{Computational Materials Science} \textbf{58}, 218--226 (2012).
doi:10.1016/j.commatsci.2012.02.005.

\bibitem{draxl2018nomad}
C. Draxl and M. Scheffler.
NOMAD: The FAIR concept for big-data-driven materials science.
\textit{MRS Bulletin} \textbf{43}, 676--682 (2018).
doi:10.1557/mrs.2018.208.

\bibitem{belsky2002icsd}
A. Belsky, M. Hellenbrandt, V. L. Karen, and P. Luksch.
New developments in the Inorganic Crystal Structure Database (ICSD): accessibility in support of materials research and design.
\textit{Acta Crystallographica Section B} \textbf{58}, 364--369 (2002).
doi:10.1107/S0108768102006948.

\bibitem{grazulis2012cod}
S. Grazulis, A. Daskevic, A. Merkys, D. Chateigner, L. Lutterotti, M. Quiros, N. R. Serebryanaya, P. Moeck, R. T. Downs, and A. Le Bail.
Crystallography Open Database (COD): an open-access collection of crystal structures and platform for world-wide collaboration.
\textit{Nucleic Acids Research} \textbf{40}, D420--D427 (2012).
doi:10.1093/nar/gkr900.

\bibitem{nimsSupercon}
National Institute for Materials Science.
SuperCon Datasheet, version v240322.
doi:10.48505/nims.4487.

\bibitem{foppiano2023supercon2}
L. Foppiano, T. C. T. Ting, J. M. Romary, and coauthors.
Automatic extraction of materials and properties from superconductors scientific literature.
\textit{Science and Technology of Advanced Materials: Methods} \textbf{3}, 2153633 (2023).
doi:10.1080/27660400.2022.2153633.

\bibitem{ishii2023ontology}
Y. Ishii and T. Sakamoto.
Structuring superconductor data with ontology: reproducing historical datasets as knowledge bases.
\textit{Science and Technology of Advanced Materials: Methods} \textbf{3}, 2223051 (2023).
doi:10.1080/27660400.2023.2223051.

\bibitem{sommer2023threedsc}
T. Sommer, D. Di Cataldo, and L. Boeri.
3DSC -- a dataset of superconductors including crystal structures.
\textit{Scientific Data} \textbf{10}, 816 (2023).
doi:10.1038/s41597-023-02721-y.

\bibitem{han2025htsc}
X.-Q. Han, Z.-F. Gao, X.-D. Wang, Z. Ouyang, P.-J. Guo, and Z.-Y. Lu.
HTSC-2025: A benchmark dataset of ambient-pressure high-temperature superconductors for AI-driven critical temperature prediction.
\textit{Chinese Physics B} \textbf{34}, 100301 (2025).
doi:10.1088/1674-1056/adf042.

\bibitem{bercx2025prxenergy}
M. Bercx, S. Ponce, Y. Zhang, G. Trezza, A. G. Ghezeljehmeidan, L. Bastonero, J. Qiao, F. O. von Rohr, G. Pizzi, E. Chiavazzo, and N. Marzari.
Charting the landscape of Bardeen-Cooper-Schrieffer superconductors in experimentally known compounds.
\textit{PRX Energy} \textbf{4}, 033012 (2025).
doi:10.1103/sb28-fjc9.

\bibitem{pires2026ambient}
P. R. Pires, T. H. B. da Silva, K. Gao, K. H. Hiorth, T. F. T. Cerqueira, T. Cavignac, P.-P. De Breuck, H.-C. Wang, D. Dangic, Y.-W. Fang, A. Sanna, W. Cui, I. Errea, P. Torma, and M. A. L. Marques.
Machine learning driven exploration of hydride superconductors at ambient pressure.
\textit{Computational Materials Today} \textbf{10}, 100052 (2026).
doi:10.1016/j.commt.2026.100052.

\bibitem{chen2026descriptors}
J. Chen, J. Peng, Y. Liang, R. Wang, H. Dong, and W. Zhang.
Interpretable descriptors enable prediction of hydrogen-based superconductors at moderate pressures.
\textit{Materials Today Physics} \textbf{63}, 102073 (2026).
doi:10.1016/j.mtphys.2026.102073.

\bibitem{wang2026hpcsd}
Z. Wang, Q. Wang, J. Duan, H. Ge, X. Luo, P. Gao, W. Zhang, J. Lv, Y. Wang, and Y. Ma.
High-Pressure Crystal Structure Database.
\textit{arXiv} 2605.14471 (2026).

\end{thebibliography}

\end{document}

\end{document}
